\documentclass[11pt,a4paper]{article}

% Packages
\usepackage[utf8]{inputenc}
\usepackage[T1]{fontenc}
\usepackage[margin=0.75in]{geometry}
\usepackage{enumitem}
\usepackage{hyperref}
\usepackage{titlesec}
\usepackage{xcolor}

% Formatting
\pagestyle{empty}
\setlength{\parindent}{0pt}
\setlist[itemize]{leftmargin=*, noitemsep, topsep=3pt}

% Section formatting
\titleformat{\section}{\Large\bfseries}{}{0em}{}[\titlerule]
\titlespacing*{\section}{0pt}{12pt}{6pt}

\titleformat{\subsection}[runin]{\bfseries}{}{0em}{}
\titlespacing*{\subsection}{0pt}{8pt}{0pt}

% Colors
\definecolor{linkcolor}{RGB}{0,102,204}
\hypersetup{
    colorlinks=true,
    urlcolor=linkcolor
}

\begin{document}

% Header
\begin{center}
    {\LARGE \textbf{ERIC MAYEUX}}\\
    \vspace{4pt}
    {\large Directeur de projets TI | IT Project Director}\\
    \vspace{4pt}
    Montréal, QC, Canada\\
    \href{mailto:mayeux.e@gmail.com}{mayeux.e@gmail.com} | 438-884-7645
\end{center}

\vspace{8pt}

\section*{Profil Professionnel}
\noindent
\textbf{Directeur de projets TI} avec 10+ ans d'expérience en \textbf{gestion de projets} d'envergure dans le secteur bancaire et financier. Expert en \textbf{mobilisation des parties prenantes}, planification et coordination de projets, et \textbf{amélioration continue}. Maîtrise des outils collaboratifs (Jira, MS Project) et des méthodologies \textbf{Agile (Scrum, Kanban)}. Certifié \textbf{PMP} (en cours) et CSM. Bilinguisme français/anglais. Fort leadership d'influence et orientation résultats.

\section*{Compétences Techniques}

\textbf{Direction de projets:} Planification et cadrage, coordination multi-équipes, livrables, échéanciers, budgets, priorités stratégiques

\textbf{Parties prenantes \& Communication:} Mobilisation, animation de réunions et ateliers, communication fluide, leadership d'influence, prise de décision

\textbf{Gestion des risques \& enjeux:} Identification, registre, plans de mitigation, suivi, amélioration continue

\textbf{Documentation:} Charte de projet, plan de projet, registre des risques, rapport d'avancement, leçons apprises

\textbf{Outils:} Jira, MS Project, Confluence, Datadog, Splunk, tableaux de bord KPI, Asana

\textbf{Méthodologies:} Agile, Scrum, SAFe, Kanban, Waterfall, PMP

\textbf{Intelligence Artificielle:} Agents IA, automatisation de gestion, RPA

\vspace{6pt}

\section*{Réalisations}
\begin{itemize}
    \item \textbf{Direction de projets} multi-équipes à la BNC : planification, coordination, livraison
    \item \textbf{Mobilisation des parties prenantes} internes et externes dans un contexte SAFe (Release Train)
    \item Mise en place de mécanismes de \textbf{suivi de l'avancement} : tableaux de bord, KPIs, rapports
    \item \textbf{Amélioration continue} : audit qualité, état des lieux, optimisation des pratiques de gestion
    \item \textbf{Animation d'ateliers} et réunions de travail pour résoudre les blocages inter-équipes
    \item Encadrement de 5+ professionnels TI dans un contexte agile international (Thaïlande)
\end{itemize}

\section*{Expérience Professionnelle}

\subsection*{Scrum Master -- Gestion de projet | Project Manager} \hfill \textit{Novembre 2025 -- Présent}\\
\textbf{Banque Nationale du Canada (BNC)} -- Montréal, QC\\
\textbf{Dirige} des projets TI stratégiques : planification, mobilisation des équipes, suivi des livrables.
\begin{itemize}
    \item \textbf{Structure les besoins} en plans de projet concrets et réalisables
    \item Encadre des équipes IT multidisciplinaires (QA, Dev, PO, Analystes) vers les objectifs
    \item Identifie et gère les \textbf{risques et enjeux} ; propose des plans de mitigation adaptés
    \item \textbf{Anime des réunions} de suivi, ateliers, comités selon les besoins du projet
    \item Produit des \textbf{rapports d'avancement} et tableaux de bord pour les parties prenantes
    \item Contribue à l'\textbf{amélioration continue} des pratiques de gestion de projet
\end{itemize}

\subsection*{Intégrateur Sénior Qualité | Senior Quality Integrator} \hfill \textit{Octobre 2023 -- Novembre 2025}\\
\textbf{Banque Nationale du Canada (BNC)} -- Montréal, QC\\
\begin{itemize}
    \item Coordination multi-sites (Thaïlande) : \textbf{gestion de plusieurs projets} simultanés
    \item Tableaux de bord (Datadog, Splunk, Jira) pour le \textbf{suivi de l'avancement}
    \item Recrutement, formation continue, \textbf{leadership d'influence} au sein de l'équipe
\end{itemize}

\subsection*{Intégrateur Qualité -- SDET | Quality Integrator -- Software Development Engineer in Test} \hfill \textit{Janvier 2022 -- Octobre 2023}\\
\textbf{Banque Nationale du Canada (BNC)} via Groupe Astek -- Montréal, QC\\
\begin{itemize}
    \item Conception et suivi de stratégie de test ; \textbf{documentation projet} et métriques qualité
    \item Stack : Selenium, AppliTools, RestAssured, JMeter, Gatling, K6
\end{itemize}

\subsection*{QA Automaticien | QA Automation Engineer} \hfill \textit{Juin 2021 -- Décembre 2021}\\
\textbf{AODocs} -- Paris, France\\
\begin{itemize}
    \item Automatisation tests non-régression ; POC Flutter ; rapports Allure
\end{itemize}

\subsection*{Automaticien CI/CD | DevOps Automation Engineer} \hfill \textit{Janvier 2020 -- Juin 2021}\\
\textbf{ENGIE DIGITAL} via Groupe Hénix -- Paris, France\\
\begin{itemize}
    \item Création CI/CD, \textbf{pilotage projets} Jira, KPIs ; scripting Python/Shell
\end{itemize}

\subsection*{Product Owner} \hfill \textit{Juin 2019 -- Septembre 2019}\\
\textbf{Hénix} -- Montrouge, France\\
\begin{itemize}
    \item Gestion du backlog, priorisation, \textbf{animation de cérémonies}, support clients SaaS
\end{itemize}

\subsection*{Développeur Logiciel | Software Developer} \hfill \textit{Mai 2017 -- Mai 2019}\\
\textbf{MEDIAPOST} via Groupe Hénix -- Paris, France\\
\begin{itemize}
    \item Développement back office, reporting, macros Excel
\end{itemize}

\subsection*{Consultant QA \& Automatisation | QA Automation Consultant} \hfill \textit{Avril 2016 -- Avril 2017}\\
\textbf{ENGIE IT} via Groupe Hénix -- Saint-Ouen, France\\
\begin{itemize}
    \item Pilotage opérationnel, coordination fournisseurs, gestion des incidents
\end{itemize}

\section*{Certifications \& Formations}

\textbf{PMP Certification} (en cours) \hfill \textit{2026}

Project Management Professional -- Project Management Institute (PMI)

\textbf{AWS Cloud Practitioner} \hfill \textit{2026}

Amazon Web Services (AWS)

\textbf{Certified Scrum Master (CSM)} \hfill \textit{2024}

Scrum Alliance

\textbf{Jira ACP-600 -- Project Administration in Jira Server} \hfill \textit{2019}

Atlassian Certified Professional

\textbf{Cursus Automaticien et DevOps} \hfill \textit{2019}\\
\textit{AFCEPF -- Montrouge, France}

\textbf{Cursus Architecture logiciel \& Analyste informaticien - Certifié Master II} \hfill \textit{2017}\\
\textit{AFCEPF-HENIX -- Malakoff, France}

\section*{Formation Académique}

\textbf{Licence en Gestion (Bachelor's Degree in Management)} \hfill \textit{2011}\\
École Supérieure de Gestion (E.S.G) -- Paris, France

\section*{Compétences Linguistiques}

\textbf{Français:} Langue maternelle (Native)\\
\textbf{Anglais:} Niveau Avancé (Advanced / Professional Working Proficiency)

\end{document}
